\subsection{20 分钟迷你版：二维、分数参数，走完前向与反向}
\phantomsection\label{app:rational-mini}
先用本例练计算链，再查下面的三维 Pre-Norm 主例。
本例固定 $X=I_2$，只监督第二位置、目标为第二类；
一个因果头的匹配宽度为 $1$、Value 宽度为 $2$，关闭 Norm 与偏置。
三维主例保留 LayerNorm，其梯度计算见第~\ref*{intro:language-model-training}~章。
取
\[
W_Q=\begin{bmatrix}0\\1\end{bmatrix},\quad W_K=\begin{bmatrix}0\\0\end{bmatrix},\quad
W_V=W_O=W_{\rm up}=W_{\rm down}=I_2,\quad
W_{\rm out}=\begin{bmatrix}1&-1\\-1/3&1/3\end{bmatrix}.
\]
\textbf{前 8 分钟：}第二行的分数均为零，因果权重 $(1/2,1/2)$；
$o=(1/2,1/2)$，第一次残差 $r=(1/2,3/2)$，
$f=\relu(rW_{\rm up})=r$，第二次残差 $h=r+fW_{\rm down}=(1,3)$。
于是 $z=hW_{\rm out}=(0,0)$、$p=(1/2,1/2)$、$L=\ln2$。

\textbf{后 12 分钟：}从 $g_z=(1/2,-1/2)$ 开始。
输出头向前传回 $g_h=g_zW_{\rm out}^{\top}=(1,-1/3)$，
两条残差路径相加得到 $g_r=2g_h=(2,-2/3)$。
依次核对全部参数梯度：
\[
\begin{aligned}
\nabla_{W_{\rm out}}L&=h^\top g_z=\begin{bmatrix}1/2&-1/2\\3/2&-3/2\end{bmatrix},\\
\nabla_{W_{\rm down}}L&=\nabla_{W_{\rm up}}L
=\begin{bmatrix}1/2&-1/6\\3/2&-1/2\end{bmatrix},\\
\nabla_{W_O}L&=\nabla_{W_V}L
=\begin{bmatrix}1&-1/3\\1&-1/3\end{bmatrix}.
\end{aligned}
\]
读取权重的梯度是 $(2,-2/3)$；经 Softmax 得到分数梯度 $(2/3,-2/3)$。
由 $q=1,k_1=k_2=0$，得到
$\nabla_{W_K}L=(2/3,-2/3)^\top$、$\nabla_{W_Q}L=(0,0)^\top$。
Query 梯度为零来自本次 $K=0$ 的参数点，不表示 Query 永远学不到。
用 $W\leftarrow W-\eta\nabla_WL$ 同步更新即可完成一步。
\texttt{python3 scripts/mini\_transformer\_demo.py} 逐一用有限差分核对全部 $24$ 个参数导数，并验证小步更新后损失下降。
\chapterpracticejump{rational-mini}
\medskip

\subsection{三维贯通实验：完整参数表与逐层中间值}
\phantomsection\label{app:transformer-intermediates}

正文\hyperref[intro:transformer-forward-pass]{贯通实验}列出了
$X,A,H,\Pi$ 四张最小矩阵。本附录给出
对应的参数与中间量，供逐项验算。

数值由 \nolinkurl{scripts/teaching_reference.py} 中的
\texttt{parameters()} 生成，展示时保留四位小数；
复算请使用未舍入的数；下文的\hyperref[app:parameter-literal]{可复制参数字典}保留全部初值。

\textbf{输入与位置。}词表大小 $5$，隐藏宽度 $3$，
输入 ID $[2,4,5]$，目标 ID $[4,5,3]$：
\[
E=\begin{bmatrix}
0&0&0\\1.2&0&-.2\\.2&-.4&.7\\-.3&1.1&.2\\.15&-.2&1.25
\end{bmatrix}\;(5\times3),\qquad
P_{\rm pos}=\begin{bmatrix}
.10&.05&-.10\\0&-.05&.12\\-.12&.10&.07
\end{bmatrix}\;(3\times3).
\]
第一步为 $X=E_{[2,4,5]}+P_{\rm pos}$，逐行 LayerNorm
采用 $\gamma=1$、$\beta=0$、$\varepsilon=10^{-5}$：
\[
U=\operatorname{LN}(X)\approx
\begin{bmatrix}
1.3833&-.4368&-.9464\\
-1.1901&1.2566&-.0665\\
-.6033&-.8061&1.4093
\end{bmatrix}.
\]

\textbf{互不相同的 Query、Key 与 Value 投影。}
\[
\begin{aligned}
W_Q&=\begin{bmatrix}.8&.2&-.1\\.1&.9&.25\\-.1&.15&.8\end{bmatrix},\\[3pt]
W_K&=\begin{bmatrix}.85&0&.2\\-.1&.9&0\\.15&.2&.75\end{bmatrix},\\[3pt]
W_V&=\begin{bmatrix}1&.2&0\\0&.9&.15\\.1&-.2&.85\end{bmatrix}.
\end{aligned}
\]
三个投影都为 $3\times3$，对应的输出是
\[
\begin{aligned}
Q&=UW_Q\approx
\begin{bmatrix}
1.1576&-.2584&-1.0047\\
-.8198&.8829&.3800\\
-.7041&-.6347&.9863
\end{bmatrix},\\[4pt]
K&=UW_K\approx
\begin{bmatrix}
1.0775&-.5824&-.4332\\
-1.1473&1.1176&-.2879\\
-.2208&-.4436&.9364
\end{bmatrix},\\[4pt]
V&=UW_V\approx
\begin{bmatrix}
1.2886&.0728&-.8700\\
-1.1968&.9062&.1320\\
-.4623&-1.1280&1.0770
\end{bmatrix}.
\end{aligned}
\]

\textbf{因果分数与残差。}先除以 $\sqrt3$，
将未来位置掩为 $-\infty$：
\[
S=QK^\top/\sqrt3+C\approx
\begin{bmatrix}
1.0583&-\infty&-\infty\\
-.9019&1.0496&-\infty\\
-.4713&-.1071&.7855
\end{bmatrix}.
\]
正文的 $A=\operatorname{softmax}(S)$
保留四位小数；随后
\[
R_{\rm raw}=AV\approx
\begin{bmatrix}
1.2886&.0728&-.8700\\
-.8876&.8025&.0074\\
-.3458&-.4345&.5215
\end{bmatrix},
\quad
W_O=\begin{bmatrix}
.9&0&.1\\.1&.85&0\\0&.15&.95
\end{bmatrix},
\]
\[
X_{\rm att}=X+R_{\rm raw}W_O\approx
\begin{bmatrix}
2.4670&-.0186&-.9976\\
-1.0186&1.7333&.2382\\
-.3247&-.3911&1.7809
\end{bmatrix}.
\]

\textbf{第二次归一化与 FFN。}网络隐藏层宽度为 $4$：
\[
W_{\rm up}=\begin{bmatrix}
.4&.2&-.3&.1\\
-.1&.5&.15&-.25\\
.2&.1&.45&.3
\end{bmatrix}\;(3\times4),
\quad
W_{\rm down}=\begin{bmatrix}
.25&-.1&.15\\
-.05&.2&.1\\
.1&.15&.15\\
.1&-.12&.18
\end{bmatrix}\;(4\times3).
\]
\[
B=\operatorname{LN}(X_{\rm att})\approx
\begin{bmatrix}
1.3601&-.3444&-1.0157\\
-1.1879&1.2585&-.0706\\
-.6739&-.7398&1.4137
\end{bmatrix},
\]
\[
F_{\rm mid}=\operatorname{ReLU}(BW_{\rm up})\approx
\begin{bmatrix}
.3753&0&0&0\\
0&.3846&.5134&0\\
.0872&0&.7274&.5417
\end{bmatrix}.
\]
$H=X_{\rm att}+F_{\rm mid}W_{\rm down}$ 为 $3\times3$，完整结果见正文第~\ref*{intro:transformer-forward-pass}~章。

\textbf{输出头与 logits。}
\[
W_{\rm out}=\begin{bmatrix}
0&.2&0&.2&-.2\\
0&.2&0&1.2&-.4\\
0&-.1&.4&.1&.2
\end{bmatrix}\;(3\times5),
\]
\[
Z=HW_{\rm out}\approx
\begin{bmatrix}
0&.5951&-.3765&.3507&-.6780\\
0&.1448&.1415&2.1027&-.4868\\
0&-.3064&.8002&-.2620&.5776
\end{bmatrix}.
\]
正文的 $\Pi=\operatorname{softmax}(Z)$ 为 $3\times5$，
目标 $[4,5,3]$ 的三行概率平均交叉熵约 $1.798430$。

\subsubsection*{输出头梯度矩阵}
第~\ref*{intro:language-model-training}~章中，$W_{\rm out}$ 的梯度形状为 $3\times5$，
十五个元素为：
\[
\frac{\partial L}{\partial W_{\rm out}}
\approx\begin{bmatrix}
0.1211&0.2472&0.1153&-0.8595&0.3760\\
0.0303&0.0404&0.1355&0.4251&-0.6313\\
0.0543&-0.0182&-0.4674&0.3902&0.0411
\end{bmatrix}.
\]


\begin{idea}
矩阵形状检查顺序：
$E_{\mathrm{ids}}(3\times3)\to X(3\times3)\to
Q,K,V(3\times3)\to A(3\times3)\to H(3\times3)
\to Z(3\times5)\to\Pi(3\times5)$。
按此顺序逐项核对行数与列数，可以检查投影、读取和输出头的矩阵乘法是否匹配。
\end{idea}

\begin{samepage}
\subsubsection*{可复制的初始参数：只有 PDF 也能运行}
\phantomsection\label{app:parameter-literal}
下面的 \texttt{parameters()} 与上面的九张参数表逐项相同，共 $99$ 个可训练数值；
键名 \texttt{P} 对应 $P_{\rm pos}$，\texttt{Wq/Wk/Wv/Wo} 对应 $W_Q/W_K/W_V/W_O$，
其余键为 \texttt{E/Wup/Wdown/Wout}。
把整个代码块复制到独立解答文件或 Notebook 单元格中，即可取得同一组初值；只依赖 NumPy。
每次调用都返回新的数组，便于从相同参数点比较两次实验。
\par\nobreak
\begin{verbatim}
import numpy as np

def parameters():
    return {
        'E': np.array([[0., 0., 0.], [1.2, 0., -.2], [.2, -.4, .7],
                       [-.3, 1.1, .2], [.15, -.2, 1.25]]),
        'P': np.array([[.1, .05, -.1], [0., -.05, .12], [-.12, .1, .07]]),
        'Wq': np.array([[.8, .2, -.1], [.1, .9, .25], [-.1, .15, .8]]),
        'Wk': np.array([[.85, 0., .2], [-.1, .9, 0.], [.15, .2, .75]]),
        'Wv': np.array([[1., .2, 0.], [0., .9, .15], [.1, -.2, .85]]),
        'Wo': np.array([[.9, 0., .1], [.1, .85, 0.], [0., .15, .95]]),
        'Wup': np.array([[.4, .2, -.3, .1], [-.1, .5, .15, -.25],
                         [.2, .1, .45, .3]]),
        'Wdown': np.array([[.25, -.1, .15], [-.05, .2, .1],
                           [.1, .15, .15], [.1, -.12, .18]]),
        'Wout': np.array([[0., .2, 0., .2, -.2],
                          [0., .2, 0., 1.2, -.4],
                          [0., -.1, .4, .1, .2]]),
    }
\end{verbatim}
在只导入 \texttt{parameters} 的参考答案中，使用本段定义代替那一行仓库导入即可；
前向与梯度计算仍由题目要求的实现完成。
\end{samepage}
